\section*{第 \ref{ch:b2:geological-hazards} 章——地质灾害与风险}

\begin{solution}{exo:b2:geological-hazards:1}
(a) $1 - e^{-30/100} = 0.259$。(b) $e^{-100/100} = e^{-1} = 0.368$：超过三分之一的
世纪未发生灾害。(c) $T = -50/\ln 0.95 = 975$ 年。
\end{solution}

\begin{solution}{exo:b2:geological-hazards:2}
(a) 城镇 A：$(1/200) \times 5000 \times 0.02 = 0.5$ 栋建筑/年；城镇 B：$(1/50) \times 800
\times 0.10 = 1.6$。(b) 城镇 B，其暴露的建筑少 6.25 倍。(c) 其灾害频率高 4 倍，
脆弱性高 5 倍，二者合计为 20 倍，
足以抵消较小的\omterm{def:b2:geological-hazards:risk}{暴露度}。
\end{solution}

\begin{solution}{exo:b2:geological-hazards:3}
(a) $\sqrt{9.81 \times 4000} = \qty{198}{m.s^{-1}}$；$\sqrt{9.81 \times 200} = \qty{44.3}{m.s^{-1}}$；
$\sqrt{9.81 \times 10} = \qty{9.90}{m.s^{-1}}$。(b) $\num{1.5e6}/198.1 = \qty{7570}{s}$，即 \qty{2.10}{h}。
(c) $800/3.6 = \qty{222.2}{m.s^{-1}}$，且 $h = c^2/g = 222.2^2/9.81 = \qty{5030}{m}$。
\end{solution}

\begin{solution}{exo:b2:geological-hazards:4}
(a) $\tan\ang{33}/\tan\ang{25} = 0.6494/0.4663 = 1.39$。(b) $F = 1$ 时 $\beta = \varphi' =
\ang{33}$。(c) 取 $c' = 0$，$F = (\cos^2\beta - r_u)\tan\varphi'/(\sin\beta\cos\beta)$，而 $\cos^2\beta
- r_u = \cos^2\beta\,(1 - \gamma_w/\gamma)$，故 $F = (1 - \gamma_w/\gamma)\tan\varphi'/\tan\beta$：
水使摩擦力减半。$(1 - 9.81/20) \times 1.393 = 0.510 \times 1.393 = 0.710$：
斜坡失稳。
\end{solution}

\begin{solution}{exo:b2:geological-hazards:5}
(a) $696/50 = 13.9$ 次/年。(b) $e^{-30/26.2} = 0.318$，与观测值 0.304 接近。
观测值略低，是因为余震的成簇分布和双震事件增加了短间隔。
(c) $1 - e^{-7/26.2} = 0.234$。
\end{solution}

\begin{solution}{exo:b2:geological-hazards:6}
(a) 震级项为 $-0.105 + 0.229 \times 0.5 = 0.0095$。对于 $d = 0$、10、30 和
\qty{100}{km}，$r = 5.57$、11.45、30.51 和 \qty{100.2}{km}，$\log_{10}\mathrm{PGA} = -0.571$、
$-0.814$、$-1.145$ 和 $-1.547$，中值 PGA 分别为 0.269、0.153、0.0715 和 \qty{0.0284}{g}。
(b) $0.153 \times 10^{0.251} = 0.153 \times 1.78 = \qty{0.272}{g}$。(c) 上下各取一个标准差：
$0.153/1.70 = \qty{0.0900}{g}$ 至 $0.153 \times 1.70 = \qty{0.260}{g}$。
\end{solution}

\begin{solution}{exo:b2:geological-hazards:7}
(a) 排序后：420（$T = 13/1 = 13$ 年）、305（6.5）、268（4.33）、251（3.25）。(b) $305 + (420 -
305)\ln(10/6.5)/\ln(13/6.5) = 305 + 115 \times 0.621 = \qty{376}{m^3.s^{-1}}$。(c) 最大洪峰
绘于 13 年处；100 年的任何取值都远超出记录范围，外推所用的极值律
其选择与参数几乎不受这 12 个数据点约束。
\end{solution}

\begin{solution}{exo:b2:geological-hazards:8}
(a) $0.6 \times 0.4 \times 0.5 \times 0.3 = 0.036$。(b) $P(\text{喷发}) = 0.6 \times 0.4 = 0.24$；
若谷中无流动，则为 $0.24 - 0.036 = 0.204$。(c) $0.6 \times 0.7 \times 0.5 \times 0.3 =
0.063$。
\end{solution}

\begin{solution}{exo:b2:geological-hazards:9}
此处 $\sin\beta\cos\beta = 0.4698$，$\cos^2\beta = 0.6710$ 及 $\tan\ang{30}/\tan\ang{35} =
0.8245$。(a) 由 $F = c'/(\gamma z\sin\beta\cos\beta) + 0.8245 = 1$ 得 $z_c = 8/(19 \times 0.4698
\times 0.1755) = \qty{5.11}{m}$。(b) $8/(19 \times 2 \times 0.4698) + 0.8245 = 0.448 + 0.825 =
1.27$。(c) $r_u = 9.81 \times 0.6710/19 = 0.3465$；$F = [8 + 38 \times (0.6710 - 0.3465) \times
0.5774]/(38 \times 0.4698) = 15.12/17.85 = 0.847$：饱和层失稳。
\end{solution}

\begin{solution}{exo:b2:geological-hazards:10}
(a) $\Pr(\text{无事件于 } [0, s+t] \mid \text{无事件于 } [0, s]) = e^{-(s+t)/T}/e^{-s/T} = e^{-t/T}$。
(b) 400 个平静年份并不改变结果：$1 - e^{-50/475} = 0.0999$，即 \qty{10}{\percent}。(c) 这
与“发生率不依赖于已流逝时间”的假设相矛盾：地震回跳后
应力需重新积累，因此刚发生地震后发生率很低，之后逐渐增大。一个震源区
包含许多处于不同循环阶段的断层；大量独立事件序列的叠加近似
服从泊松过程，而对每一序列而言，已知的只是其发生率。
\end{solution}

\begin{solution}{exo:b2:geological-hazards:11}
(a) $A = 0.5 \times (4000/10)^{1/4} = 0.5 \times 4.47 = \qty{2.24}{m}$。(b) $\lambda = cT$：
深海中为 $198.1 \times 1200 = \qty{238}{km}$，
水深 \qty{10}{m} 处为 $9.90 \times 1200 = \qty{11.9}{km}$。(c) 振幅相对水深已不再很小（\qty{2.2}{m}，
水深 \qty{10}{m}），因此波浪变陡并可能破碎；部分能量被
海底与海湾的形状所反射或聚焦；陆地上的爬高取决于地面的坡度与
粗糙度。
\end{solution}

\begin{solution}{exo:b2:geological-hazards:12}
(a) $10/5.8 + 4 = 1.72 + 4 = \qty{5.72}{s}$。(b) S 波到达 \qty{50}{km} 处需 $50/3.2 =
\qty{15.6}{s}$：$15.6 - 5.72 = \qty{9.90}{s}$ 的预警时间；在 \qty{100}{km} 处，$31.3 - 5.72 =
\qty{25.5}{s}$。(c) $3.2 \times 5.72 = \qty{18.3}{km}$。
\end{solution}

\begin{solution}{pb:b2:geological-hazards:1}
中间步骤中的对数保留四位小数。
\textbf{1.} $10^{4-5} = 0.1$、$10^{4-6} = 0.01$ 和 $10^{4-7} = 0.001$ 每年。
\textbf{2.} 10、100 和 1000 年。
\textbf{3.} $10^{4-7.5} = \num{3.16e-4}$ 每年；$0.001 - \num{3.16e-4} = \num{6.84e-4}$ 每年。
\textbf{4.} $0.1 - 0.01 = 0.09$ 和 $0.01 - 0.001 = 0.009$ 每年。
\textbf{5.} $1 - e^{-0.01 \times 50} = 0.393$。
\textbf{6.} $1 - e^{-0.001 \times 50} = 0.0488$。
\textbf{7.} $r = \sqrt{10^2 + 5.57^2} = \qty{11.45}{km}$；$\log_{10} r = 1.0588$，$-0.778 \times
1.0588 = -0.8237$。
\textbf{8.} $\mu = -0.105 + 0.229(M - 6) - 0.8237$：$-1.0432$、$-0.8142$ 和 $-0.6424$；中值
为 0.0905、0.153 和 \qty{0.228}{g}。
\textbf{9.} $10^{-0.8142 + 0.230} = \qty{0.260}{g}$。
\textbf{10.} $z = (\log_{10} 0.3 + 0.8142)/0.230 = (-0.5229 + 0.8142)/0.230 = 1.266$，且 $1 -
\Phi(1.266) = 0.103$。
\textbf{11.} $-0.105 + 0.229 \times 1.5 - 0.8237 = -0.5852$：\qty{0.260}{g}。
\textbf{12.} 超过 \qty{0.2}{g} 的概率（$\log_{10} 0.2 = -0.6990$）分别为 0.0672
（$z = 1.497$）、0.308（$z = 0.501$）和 0.597（$z = -0.246$）：$\lambda = 0.09 \times 0.0672 + 0.009
\times 0.308 + \num{6.84e-4} \times 0.597 = \num{6.05e-3} + \num{2.77e-3} + \num{4.08e-4} =
\num{9.23e-3}$ 每年。
\textbf{13.} 对于 \qty{0.3}{g}：$z = 2.262$、1.266 和 0.520，概率分别为 0.0118、0.103 和
0.302；$\lambda = \num{1.07e-3} + \num{9.24e-4} + \num{2.06e-4} = \num{2.20e-3}$ 每年。
\textbf{14.} 对于 \qty{0.4}{g}（$\log_{10} 0.4 = -0.3979$）：$z = 2.805$、1.810 和 1.063，
概率分别为 0.00251、0.0352 和 0.144；$\lambda = \num{2.26e-4} + \num{3.17e-4} + \num{9.84e-5} =
\num{6.41e-4}$ 每年。
\textbf{15.} $-\ln 0.9/50 = \num{2.107e-3}$ 每年，即\omterm{def:b2:geological-hazards:return}{重现期}为 475 年。
\textbf{16.} 介于 \qty{0.3}{g}（\num{2.20e-3}）与 \qty{0.4}{g}（\num{6.41e-4}）之间。
\textbf{17.} 斜率为 $\ln(6.41/22.0)/\ln(4/3) = -1.233/0.2877 = -4.29$，且 $a = 0.3 \times
(2.107/2.20)^{-1/4.29} = 0.3 \times 1.0101 = \qty{0.303}{g}$。
\textbf{18.} 在 \qty{0.3}{g} 处，5--6 级区间贡献 \qty{48.5}{\percent} 的发生率，6--7 级
贡献 \qty{42.1}{\percent}，7--7.5 级仅贡献 \qty{9.4}{\percent}。即使是最大的地震，其中值也仅为
\qty{0.260}{g}，低于 \qty{0.3}{g}：该水平仅在散布的上尾部分才能达到，而
大量中等地震提供了其中的大部分。
\textbf{19.} $4/(19 \times 3 \times 0.4330) + \tan\ang{34}/\tan\ang{30} = 0.162 + 1.168 = 1.33$。
\textbf{20.} $r_u = 9.81 \times 0.75/19 = 0.387$；$F = [4 + 57 \times (0.75 - 0.387) \times
0.6745]/(57 \times 0.4330) = 17.95/24.68 = 0.727$：饱和斜坡失稳。
\textbf{21.} 设地下水位位于层厚的分数 $m$ 处时，$r_u = 0.387\,m$。$F = 1$ 时
$0.75 - 0.387\,m = (24.68 - 4)/(57 \times 0.6745) = 0.538$：$m = 0.548$，即地下水位位于层厚的
\qty{55}{\percent} 处。
\textbf{22.} 将 $F = 1$ 除以 $\gamma z$：$0.0702 + (0.75 - 0.4330\,k) \times 0.6745 = 0.4330 +
0.75\,k$，故 $k_c = (0.0702 + 0.5059 - 0.4330)/(0.2921 + 0.75) = 0.1431/1.042 = 0.137$。
\textbf{23.} $\log_{10} 0.137 = -0.8633 = -0.105 + 0.229(M - 6) - 0.8237$ 给出 $M = 6.29$。此类
地震的发生率为 $10^{4 - 6.29} - 10^{-3.5} = \num{4.81e-3}$ 每年，约每 208 年一次，
而问题~17 中的 475 年 PGA 超过临界值 \qty{0.137}{g} 一倍以上：按此分析，
干燥斜坡在远低于 475 年地震动的强度下即失稳。
\textbf{24.} 在该镇，50 年内\omterm{def:b2:geological-hazards:poisson}{超越概率}为 \qty{10}{\percent} 的\omterm{def:b2:geological-hazards:psha}{峰值地面加速度}
为 \textbf{\qty[text-series-to-math=true]{0.303}{g}}。
\end{solution}
